% =====================================================================
\section{Training and Results}
\label{sec:results}
% =====================================================================

This section reports how the rhizome was trained and how it compares with the best models of the existing \texttt{fno-21cm} pipeline on the task of predicting the neutral-fraction field $\xHI$, and later also the brightness temperature $T_b$, from the matter density of \cmfast{} lightcones. The experiments followed the order of \cref{tab:rhizome-variants}: 2-D slices, native 3-D windows, multi-field 3-D, and finally slice-wise reconstruction of full cones. \Cref{fig:metric-ratios} summarizes the outcome.

\begin{figure}[!htbp]
  \centering
  \figplaceholder[\linewidth]{metric_ratios.pdf}{Headline metrics relative to fno.}
  \caption{Headline results on 199 held-out multi-field test cones. Each bar is the error of a rhizome model divided by the error of the best \texttt{fno-21cm} multi-field model (\texttt{tbs+histtb}); values below the dashed line are better than the baseline. Front width is measured as the excess over the true front width of 2.77\,Mpc. The reported slice-wise errors are lower on the listed metrics except the small-scale $T_b$ power before reionization, which a power term in the loss reduces (\cref{sec:results-absorption}). Numbers in \cref{tab:main-results}.}
  \label{fig:metric-ratios}
\end{figure}

% ---------------------------------------------------------------------
\subsection{Datasets}
\label{sec:datasets}

All data come from one suite of 6,600 \cmfast{} \citep{mesinger2011cmfast} simulations (version 4.1.1), each with eleven sampled simulation parameters. Every simulation provides a 3-D lightcone of $140\times140$ transverse cells over a 200\,Mpc box (1.43\,Mpc per cell). The native line-of-sight length depends on cosmology and is between 2,074 and 2,919 cells at the same spacing. The lightcones contain the density, the neutral fraction, the brightness temperature and the LOS velocity, among other fields. Two properties of the data matter for the models:
\begin{itemize}
  \item Each lightcone reuses the periodic 200\,Mpc spatial realization along the LOS, with a tiling period of 140 cells. The fields can evolve with redshift along that direction, so this is repeated spatial sampling rather than an identical copy of the evolving field.
  \item Ionized cells have $\xHI$ exactly zero; partially ionized values appear mostly where the lightcone interpolates between two simulation snapshots.
\end{itemize}
33 simulations (0.5\%) contain non-finite $T_b$ values and are excluded from every multi-field split. \Cref{tab:datasets} lists the splits used.

\begin{table*}[!tp]
  \centering
  \caption{Training, validation and test splits. All test sets are held out by simulation (cone), never by slice.}
  \label{tab:datasets}
  \small
  \begin{tabular}{@{}p{2.9cm}p{4.3cm}p{2.6cm}p{4.2cm}@{}}
    \toprule
    Split & Used for & Train / val / test cones & Samples \\
    \midrule
    2-D slice cache & 2-D sweep and first 2-D model & 5,280 train (own split) & 16 window + 2 background slices per cone \\
    $\xHI$ (fno split) & 3-D rhizome; fair slice-wise model & 1,600 / 200 / 200 & 3-D windows; or 48 + 6 slices per cone \\
    Multi-field & 3-D multi-field; slice-wise multi-field & 1,594 / 199 / 199 & as above, cleaned $T_b$ \\
    Multi-field, all cones & best slice-wise multi-field model & 6,169 / 199 / 199 & 16 + 2 slices per cone (98.5\,k) \\
    \bottomrule
  \end{tabular}
\end{table*}

The $\xHI$ and multi-field splits are those of the \texttt{fno-21cm} pipeline. All comparisons with \texttt{fno} models below use identical test cones and the same evaluation code.

% ---------------------------------------------------------------------
\subsection{Training protocol}
\label{sec:training}

\paragraph{Losses.} $\xHI$ is trained with binary cross-entropy (BCE) on the logit. In the multi-field models, $T_b$ is added as a mean-squared error on the normalized field, through the structured head of \cref{eq:structured-head}. The choice of BCE turned out to matter beyond the operator (\cref{sec:results-3d-mf}). The \texttt{fno} baselines train $\xHI$ with a mean-squared error through a sigmoid, whose gradient vanishes when the prediction saturates.

\paragraph{Optimization and stopping.} The learning rate is $2\times10^{-3}$ throughout; the 3-D models use warmup and cosine decay. The 2-D models are stopped by a predefined plateau rule on the validation loss, and the checkpoint with the best validation score is kept. The 3-D models are trained for 20 epochs of 6,400 windows and then continued with a short \emph{warm restart} (10 epochs at $5\times10^{-4}$ with cosine decay), which improved test RMSE by 3--6\%, against about 1\% for a continuation at constant small learning rate. The 2-D models use the dihedral symmetries of the square slice as data augmentation.

\paragraph{Meaning of the neutral-fraction loss.} BCE accepts targets in $[0,1]$, including partially ionized simulation cells; it need not mean that the output is a Bernoulli sample. For deterministic targets it acts as a bounded-fraction regression loss. If the compressed inputs leave some morphology unresolved, both BCE and squared error are minimized pointwise by the conditional mean of the target. BCE alone therefore does not eliminate averaging of uncertain boundaries or produce a distribution of possible fields.

\paragraph{Hyperparameters.} \Cref{tab:hyperparameters} lists the main settings. Every model is a single training run with seed 0, except for one repeat of the main multi-field configuration with seed 1 (\cref{sec:results-variations}).

\paragraph{Recorded training details of the slice-wise multi-field models.} The run metadata fix the following settings. The optimizer is Adam without weight decay or gradient clipping. Validation runs every 1,000 steps. The learning rate is halved after four validations without a relative improvement of 0.2\%, down to $2\times10^{-5}$; training stops after eight such validations at the minimum rate, but not before 20,000 steps. The checkpoint with the lowest validation score, the mean of the normalized $\xHI$ and $T_b$ mean-squared errors, is kept. The loss is BCE on $\xHI$ plus the normalized $T_b$ MSE with weight 1. Density is divided by 10, $\xHI$ is not rescaled, and $T_b$ is normalized as $(T_b+28.1\,\mathrm{mK})/56.1\,\mathrm{mK}$ with training-set statistics. Per cone, 16 slices are drawn stratified in the global history within $0.02\le\bar x_{\mathrm{HI}}\le0.98$, plus two background slices outside that window, over $5.0\le z\le24.97$. The all-cones cache gives 98,509 training and 3,080 validation slices; the slice-level test set has 3,129 slices from 199 cones.

\begin{table*}[!tp]
  \centering
  \caption{Main hyperparameters.}
  \label{tab:hyperparameters}
  \small
  \begin{tabular}{@{}p{3.6cm}p{5.0cm}p{5.4cm}@{}}
    \toprule
    & 2-D slice / slice-wise models & 3-D window model \\
    \midrule
    width $\dv$ & 64 (128 in early runs) & 48 \\
    Fourier modes $\kmax$ & 24 & (36, 36, 24), rank-48 factorized \\
    updates $\Niter$ / step $\Delta t$ & 6 / 0.5 & 6 / 0.5 \\
    batch / sample & 32 slices & native window, $140\times140\times256$, including $2\times32$ halo \\
    learning rate & $2\times10^{-3}$, plateau-stopped & $2\times10^{-3}$ (warmup + cosine), 20 epochs + warm restart \\
    precision & float32 & bfloat16 autocast, float32 FFT \\
    parameters & grow as $\dv^2$: 1.1\,M (w32) to 18.3\,M (w128) & 5.8\,M \\
    training time & 2.8--3.5\,h (A100) & $\sim$12\,h + 6\,h restart (H200) \\
    \bottomrule
  \end{tabular}
\end{table*}

% ---------------------------------------------------------------------
\subsection{Evaluation metrics}
\label{sec:metrics}

Per-voxel RMSE is reported, but it is a weak guide to physical quality: a model can lower its RMSE by blurring fronts. The evaluation therefore uses the physical suite of the \texttt{fno-21cm} pipeline, computed on full test cones. Most metrics are evaluated in the \emph{active} stage of reionization, slices with mean neutral fraction between 0.05 and 0.95, and some also \emph{before reionization}, $\bar x_{\mathrm{HI}}\ge0.98$, where $T_b$ is in absorption. \Cref{tab:metrics} lists them.

\begin{table*}[!tp]
  \centering
  \caption{Evaluation metrics. All are computed per test cone and summarized by the median over cones.}
  \label{tab:metrics}
  \small
  \begin{tabular}{@{}p{4.3cm}p{9.6cm}@{}}
    \toprule
    Metric & What it measures \\
    \midrule
    $|\Delta\tau_e|/\sigma_{\mathrm{Planck}}$ & error in the CMB optical depth implied by the predicted global history, in units of the Planck uncertainty \\
    $|\Delta z|$ at $\bar x=0.25/0.5/0.75$; duration & timing of reionization: error in the redshift at which the global neutral fraction crosses each value, and in the duration \\
    $\xHI$ PDF $W_1$ & Wasserstein-1 distance between predicted and true one-point distributions of $\xHI$ \\
    partially ionized fraction & fraction of cells with intermediate $\xHI$ (truth $\approx0.09$); measures hedging \\
    $b_1$ pull; $P_\varepsilon$ ratio & EFT diagnostics: error in the fitted linear bias in units of its uncertainty; ratio of stochastic power to the truth (ideal 1) \\
    bubble-size $W_1$ & Wasserstein-1 distance (Mpc) between mean-free-path bubble-size distributions \\
    power error (large / intermediate / small scales) & $|P_{\mathrm{pred}}/P_{\mathrm{true}}-1|$ of the transverse $\xHI$ power spectrum in three $k$ ranges; small scales mean $k>1\,\mathrm{Mpc}^{-1}$ \\
    $1-r$, high $k$ & small-scale incoherence: one minus the cross-correlation \cref{eq:cross-correlation}, averaged over $k>1\,\mathrm{Mpc}^{-1}$ \\
    edge $L^2$ / $H^1$, $\pm2$\,Mpc & error (and error in gradients) in a $\pm2$\,Mpc band around true ionization fronts \\
    front width & width of the predicted ionization fronts (Mpc), compared with the true width \\
    $T_b$ PDF $W_1$; $T_b$ power deviation & distance between $T_b$ distributions (mK); power difference in units of the single-slice sample variance ($\sigma$) \\
    \bottomrule
  \end{tabular}
\end{table*}

The cross-correlation coefficient between prediction and truth,
\begin{equation}
  \label{eq:cross-correlation}
  r(k) = \frac{P_{\mathrm{pred}\times\mathrm{true}}(k)}
              {\sqrt{P_{\mathrm{pred}}(k)\,P_{\mathrm{true}}(k)}},
\end{equation}
deserves special attention. The power ratio measures whether a model produces the right \emph{amount} of structure at each scale; $r(k)$ measures whether that structure is in the right \emph{place}. In the earlier 3-D operator study of the \texttt{fno-21cm} pipeline, small-scale coherence predicted the bias in the bubble-size distribution across architectures (Spearman $\rho=0.88$), while small-scale power amplitude did not. Getting fronts in the right place matters more for bubble statistics than getting the right amount of small-scale detail.

% ---------------------------------------------------------------------
\subsection{Baselines}
\label{sec:baselines}

All baselines come from the \texttt{fno-21cm} pipeline and were trained there on the same splits:
\begin{itemize}
  \item \textbf{\texttt{cnn\_whno\_es}} (3.5\,M parameters): the \texttt{fno} 3-D $\xHI$ model trained on native LOS windows. It has the U-FNO layout, a convolutional local path next to a learned global path, at a size comparable to the rhizome. Trained for 30 epochs (44.7\,h on an A100).
  \item \textbf{\texttt{es+histtb}} and \textbf{\texttt{tbs+histtb}}: the corresponding multi-field models, which receive density, velocity and both emulated global histories. \texttt{tbs+histtb} additionally uses the structured $T_b$ head and is the strongest \texttt{fno} model. It corresponds to rhizome variant C below and is labelled ``UFNO'' in the head-to-head figures.
  \item \textbf{2-D \texttt{fno} models} (LWF, WHNO, LocalFNO, U-FNO), which predict one slice from that slice alone.
\end{itemize}
The 203\,M-parameter implementation of U-FNO \citep{wen2022ufno} used in the earlier operator study (\cref{sec:ufno}) is not used as a baseline here, because its capacity is not comparable.

% ---------------------------------------------------------------------
\subsection{2-D slices: line-of-sight context dominates}
\label{sec:results-2d}

A sweep around a centre configuration (width 64, 24 modes, six updates) showed that the largest single factor is not a hyperparameter of the operator but the LOS context in the input (\cref{fig:los-context}). With only the target slice, the rhizome reaches a mixed-slice RMSE of 0.174. Adding the six neighbouring slices brings it to 0.138, and adding the band means out to $\pm32$ cells to 0.112, a 36\% reduction in total. By comparison, the other settings matter little: width 128 is 3\% better than width 64 at four times the parameters, 12 or 48 modes are worse than 24, and three or ten updates are slightly worse than six in this short sweep. In full-length training of the multi-field model, 3 and 10 updates match 6 within the seed-to-seed spread (\cref{sec:results-variations}). At equal compute, more training slices beat more width.

\begin{figure}[!htbp]
  \centering
  \figplaceholder[\linewidth]{los_context.pdf}{LOS context chart.}
  \caption{Effect of line-of-sight context on the 2-D rhizome (sweep, 15\,k steps, one seed). Mixed slices have mean $\xHI$ between 0.05 and 0.95.}
  \label{fig:los-context}
\end{figure}

\Cref{tab:results-2d} compares the 2-D rhizome with the 2-D \texttt{fno} models on the same 7,036 validation slices. The rhizome reduces the error of the next best model by about 39\%, and its ionization fronts are five times sharper. Most of this gap, however, is LOS information: the \texttt{fno} models see a single slice, and the single-slice rhizome (0.174) performs at their level. The 2-D comparison therefore shows what the input buys rather than what the operator buys. The fully trained 2-D model (width 128, plateau-stopped at 68\,k steps) reached a test mixed-slice RMSE of 0.099 and an ionized-region IoU of 0.914.

\begin{table}[!htbp]
  \centering
  \caption{2-D rhizome against 2-D \texttt{fno} models on 7,036 validation slices (rhizome: sweep run at 15\,k steps). Only the rhizome sees LOS context. RMSE: slice RMSE; IoU: ionized-region intersection over union; front: front width (truth 3.0\,Mpc).}
  \label{tab:results-2d}
  \small
  \begin{tabular}{@{}lrrr@{}}
    \toprule
    Model & RMSE & IoU & Front (Mpc) \\
    \midrule
    \textbf{Rhizome} (LOS bands) & \textbf{0.109} & \textbf{0.873} & \textbf{6.1} \\
    \texttt{fno} LWF & 0.178 & 0.755 & 32.6 \\
    \texttt{fno} WHNO & 0.181 & 0.755 & 31.0 \\
    \texttt{fno} LocalFNO & 0.184 & 0.742 & -- \\
    \texttt{fno} U-FNO & 0.195 & 0.722 & -- \\
    \bottomrule
  \end{tabular}
\end{table}

% ---------------------------------------------------------------------
\subsection{3-D windows: \texorpdfstring{$\xHI$}{xHI}}
\label{sec:results-3d}

The 3-D rhizome sees the same information as the 3-D \texttt{fno} model: native LOS windows of the density field. This removes the LOS-input difference, although capacity, loss and optimization still differ, so it does not isolate the interaction law alone. On the 200 test cones (\cref{tab:results-3d}), the rhizome at width 48 reaches an RMSE of 0.0594 after the warm restart, against 0.0684 for \texttt{cnn\_whno\_es}, with 5.8\,M against 3.5\,M parameters. Narrower was better (width 64 gave 0.0628), and a learning rate of $3\times10^{-4}$ was 13\% worse than $2\times10^{-3}$.

\begin{table*}[!tp]
  \centering
  \caption{3-D $\xHI$ on 200 test cones. \emph{Left:} test RMSE. \emph{Right:} physical suite (medians over cones), rhizome w48 after the warm restart.}
  \label{tab:results-3d}
  \small
  \begin{tabular}[t]{lrr}
    \toprule
    Model & Params & RMSE \\
    \midrule
    rhizome w48, 20 ep. & 5.8\,M & 0.0619 \\
    \quad + warm restart & 5.8\,M & \textbf{0.0594} \\
    rhizome w64, + restart & -- & 0.0628 \\
    \texttt{cnn\_whno\_es}, 30 ep. & 3.5\,M & 0.0684 \\
    \bottomrule
  \end{tabular}
  \hspace{1.2em}
  \begin{tabular}[t]{lrr}
    \toprule
    Metric & Rhizome & \texttt{fno} \\
    \midrule
    $|\Delta\tau_e|/\sigma_{\mathrm{Planck}}$ & \textbf{0.019} & 0.058 \\
    $|\Delta z(\bar x=0.5)|$ & \textbf{0.028} & 0.098 \\
    $b_1$ pull & \textbf{1.58} & 3.41 \\
    $P_\varepsilon$ ratio (ideal 1) & \textbf{0.89} & 0.77 \\
    bubble-size $W_1$ (Mpc) & \textbf{0.81} & 1.39 \\
    small-scale power error & \textbf{0.222} & 0.336 \\
    edge $H^1$, $\pm2$\,Mpc & \textbf{0.137} & 0.145 \\
    front width (Mpc) & \textbf{5.9} & 7.5 \\
    \bottomrule
  \end{tabular}
\end{table*}

The margin is much larger on the physical metrics than on RMSE. The rhizome is three times more accurate in $\tau_e$ and in the reionization midpoint, halves the bias error, and comes 40\% closer in bubble sizes. Its small-scale power and coherence also hold up better at every stage of reionization.

% ---------------------------------------------------------------------
\subsection{3-D multi-field: inputs, heads and losses}
\label{sec:results-3d-mf}

The multi-field models predict $\xHI$ and $T_b$ from density and LOS velocity. Four 3-D rhizome configurations were trained to separate the contributions of inputs and output head (\cref{tab:results-3d-mf}).

\begin{table}[!htbp]
  \centering
  \caption{3-D multi-field ablation, test RMSE on 199 cones ($\xHI$ / $T_b$ in mK). History channels are the band means of the emulated global $\xHI(z)$ and $T_b(z)$.}
  \label{tab:results-3d-mf}
  \footnotesize
  \setlength{\tabcolsep}{3pt}
  \begin{tabular}{@{}lp{2.6cm}rr@{}}
    \toprule
    Run & Inputs / head & 20 epochs & + restart \\
    \midrule
    A0 & density $\to\xHI$ only & 0.0602 & 0.0567 \\
    A & + velocity, + $T_b$ (plain head) & 0.0648 / 9.49 & 0.0622 / 8.71 \\
    B & A + history channels & 0.0509 / 6.26 & 0.0494 / 6.07 \\
    \textbf{C} & B + structured $T_b$ head & 0.0507 / 5.97 & \textbf{0.0495 / 5.79} \\
    \midrule
    \texttt{fno} \texttt{es+histtb} & = B & 0.0577 / 6.63 & \\
    \texttt{fno} \texttt{tbs+histtb} & = C & 0.0560 / 6.54 & \\
    \bottomrule
  \end{tabular}
\end{table}

\begin{itemize}
  \item \textbf{History channels give the largest improvement in this ablation.} Going from A to B reduces the $\xHI$ error by 21\% and the $T_b$ error by 34\%.
  \item \textbf{Velocity and a $T_b$ target alone do not help $\xHI$}; A is even slightly worse than A0 (one seed each).
  \item \textbf{The structured head mainly helps $T_b$}: without it, $T_b$ power is about 10\% low at every stage. It also improves the reionization timing by 15--30\%.
  \item \textbf{The density-only rhizome A0 already beats \texttt{fno}~\texttt{es+histtb}} (0.0567 vs 0.0577), although the latter also receives velocity and both histories.
\end{itemize}
On the physical suite, model C at 20 epochs, before the warm restart, beats \texttt{tbs+histtb} by a factor of 5 in $\tau_e$ (0.014 vs 0.076\,$\sigma$), and is closer in $b_1$ (pull 1.52 vs 2.28), $P_\varepsilon$ (0.935 vs 0.793), bubble sizes ($W_1$ 0.71 vs 1.21\,Mpc) and front width (4.53 vs 6.17\,Mpc; truth 2.77). Before reionization, \texttt{fno} adds spurious small-scale $T_b$ structure ($+12.8\sigma$ against $+3.9\sigma$). After the warm restart, model C improves further; that checkpoint is used in \cref{tab:main-results} ($\tau_e$ 0.012\,$\sigma$, $b_1$ pull 1.35, $P_\varepsilon$ 0.945, $W_1$ 0.69\,Mpc, front width 4.38\,Mpc).

\paragraph{Influence of the loss on the performance gap.} Both model families receive the same emulated global history, yet the global $\bar x_{\mathrm{HI}}(z)$ of the \texttt{fno} models sits 0.002--0.003 too low from $z\approx8$ to 25. In truly neutral cells at $z>15$ the mean error is $-0.0020$ and $-0.0030$ for the two \texttt{fno} runs against $-0.00006$ for the rhizome. A plausible contributor is the loss. For a neutral-fraction logit $\ell$, prediction $p=\mathrm{sigmoid}(\ell)$ and target $y$, $\partial_\ell\mathcal L_{\mathrm{BCE}}=p-y$, whereas a squared-error loss gives $2(p-y)p(1-p)$. Both approach zero at a correct saturated prediction, but squared error has an additional sigmoid factor that slows correction near the endpoints. This is consistent with the observed neutral-cell offset; a loss-matched retraining is needed to establish its contribution. This uniform offset integrates into a $\tau_e$ bias. A fair comparison of $\tau_e$ and the global history therefore needs the \texttt{fno} models retrained with BCE. The offset alone does not establish an explanation for the morphology differences; the effect of changing the loss on those metrics also needs to be measured.

% ---------------------------------------------------------------------
\subsection{Slice-wise reconstruction and native 3-D models}
\label{sec:results-slicewise}

The slice-wise variant (\cref{fig:slicewise}) applies the 2-D model to every slice of a cone. For a fair comparison it was retrained on a new slice cache built from the same 1,600 training cones as the 3-D models, with otherwise unchanged settings. \Cref{tab:results-slicewise} compares it with the 3-D rhizome and with \texttt{fno} on all 200 test cones.

\begin{table*}[!tp]
  \centering
  \caption{Slice-wise 2-D reconstruction against the 3-D models, $\xHI$, 200 test cones, all trained on the same 1,600 cones.}
  \label{tab:results-slicewise}
  \small
  \begin{tabular}{@{}lrrrr@{}}
    \toprule
    Metric & 2-D w128 & \textbf{2-D w64} & 3-D rhizome w48 & \texttt{cnn\_whno\_es} \\
    \midrule
    $|\Delta\tau_e|/\sigma_{\mathrm{Planck}}$ & \textbf{0.011} & \textbf{0.011} & 0.019 & 0.058 \\
    $|\Delta z(\bar x=0.5)|$ & \textbf{0.011} & 0.014 & 0.028 & 0.098 \\
    $\xHI$ PDF $W_1$, active & \textbf{0.011} & 0.013 & 0.021 & 0.030 \\
    $b_1$ pull & 1.01 & \textbf{0.85} & 1.58 & 3.41 \\
    $P_\varepsilon$ ratio (ideal 1) & \textbf{0.95} & 0.94 & 0.89 & 0.77 \\
    bubble-size $W_1$ (Mpc) & 0.59 & \textbf{0.56} & 0.81 & 1.39 \\
    power error, mid $k$ & \textbf{0.062} & 0.073 & 0.174 & 0.172 \\
    $1-r$, high $k$ & \textbf{0.119} & 0.120 & 0.199 & 0.181 \\
    edge $H^1$, $\pm2$\,Mpc & \textbf{0.109} & 0.110 & 0.138 & 0.145 \\
    front width (Mpc; truth 2.65) & \textbf{4.03} & 4.10 & 5.91 & 7.54 \\
    \bottomrule
  \end{tabular}
\end{table*}

The slice-wise model wins on nearly every physical metric, trains about seven times faster than the 3-D rhizome, and has no jitter between slices: the roughness from slice to slice matches the truth to 2--3\%, while both 3-D models are slightly over-smoothed along the LOS. Its advantage comes from front \emph{placement}: small-scale incoherence drops from 0.20 to 0.12, the mid-scale power error by 60\%, and the edge error by 21\% (significant in a paired bootstrap over cones: 2,000 resamples of the per-cone differences, 95\% percentile interval). The edge metrics use 175 of the 200 test cones: in the other 25, which reionize late, no cell falls below $\xHI=0.5$ before $z=5$, so there is no front to measure. Width barely matters, so width 64 (2.8\,h of training against 5.0\,h) is the choice. Per cone, the 2-D model wins 11 of 13 inspected cones. It loses on the earliest reionizer, cone 27 ($z_{\mathrm{mid}}\approx15.9$), where the longer context of the 3-D model helps slightly.

The error maps in \cref{fig:slicewise-errors} show the qualitative difference. The slice-wise errors are thin rims along bubble walls, whereas the 3-D models, \texttt{fno} strongly and the 3-D rhizome weakly, miss whole ionized regions, and \texttt{fno} leaves a faint ionized haze over neutral gas.

\begin{figure}[!htbp]
  \centering
  \figplaceholder[0.9\linewidth]{rhizome/slicewise_xhi_error_maps.png}{Slice-wise error maps.}
  \caption{Signed $\xHI$ error at $\bar x_{\mathrm{HI}}\approx0.4$ on eight test cones (red: predicted too neutral, blue: too ionized). Columns: truth, slice-wise 2-D rhizome (w64), 3-D rhizome (w48, warm restart), \texttt{fno} \texttt{cnn\_whno\_es}; per-panel RMSE in the corner. The bottom row is cone 27, the earliest reionizer, where all models predict ionized gas as too neutral.}
  \label{fig:slicewise-errors}
\end{figure}

For this simulation suite, a 2-D model with $\pm32$ cells ($\approx\pm46$\,Mpc) of LOS context is more accurate and cheaper than the tested 3-D windows. This does not establish a universal physical interaction length. The periodic tiling, compressed context channels, padding restrictions and optimization budget may all contribute to the result.

% ---------------------------------------------------------------------
\subsection{Multi-field slice-wise results}
\label{sec:main-results}

The final model combines the slice-wise approach with the multi-field inputs and the structured head of model C: 13 LOS bands each of density and velocity, global-history band means as scalars, and $\xHI$ and $T_b$ as outputs (\cref{sec:rhizome-variants}). It was trained on all 6,169 usable training cones (3.5\,h on an A100) and evaluated on the 199 multi-field test cones. \Cref{tab:main-results} compares it with 3-D rhizome C and with the best \texttt{fno} model.

\begin{table*}[!tp]
  \centering
  \caption{Main results, 199 multi-field test cones (medians). ``Slice-wise'' is the multi-field slice-wise 2-D rhizome trained on all cones; ``+ $T_b$ power'' adds the power term of \cref{sec:results-absorption}. 3-D rhizome C is the checkpoint after the warm restart. Best value per row in bold.}
  \label{tab:main-results}
  \small
  \setlength{\tabcolsep}{4pt}
  \begin{tabular}{@{}lrrrr@{}}
    \toprule
    Metric & Slice-wise & + $T_b$ power & 3-D rhizome C & \texttt{fno} \texttt{tbs+histtb} \\
    \midrule
    $|\Delta\tau_e|/\sigma_{\mathrm{Planck}}$ & 0.010 & \textbf{0.008} & 0.012 & 0.076 \\
    $|\Delta z(\bar x=0.5)|$ / duration & \textbf{0.010 / 0.018} & \textbf{0.010 / 0.018} & 0.011 / \textbf{0.018} & 0.017 / 0.044 \\
    $\xHI$ PDF $W_1$, active & \textbf{0.0095} & 0.0098 & 0.0156 & 0.0271 \\
    $b_1$ pull & 0.99 & \textbf{0.93} & 1.35 & 2.28 \\
    $P_\varepsilon$ ratio (ideal 1) & \textbf{0.954} & -- & 0.945 & 0.793 \\
    bubble-size $W_1$, active (Mpc) & 0.56 & \textbf{0.48} & 0.69 & 1.21 \\
    $\xHI$ power error, high $k$ & 0.168 & \textbf{0.163} & 0.223 & 0.311 \\
    $1-r$, high $k$ & \textbf{0.096} & 0.102 & 0.151 & 0.174 \\
    edge $H^1$, $\pm2$\,Mpc & \textbf{0.106} & 0.112 & 0.131 & 0.147 \\
    front width (Mpc; truth 2.77) & \textbf{3.89} & 3.94 & 4.38 & 6.17 \\
    $T_b$ PDF $W_1$, active (mK) & 0.79 & \textbf{0.77} & 1.06 & 1.19 \\
    pre-reion.\ $T_b$ power dev., $k=0.46$ & $+24.8\sigma$ & $\bm{+1.5\sigma}$ & $+4.1\sigma$ & $+12.8\sigma$ \\
    \bottomrule
  \end{tabular}
\end{table*}

Against \texttt{fno}, the slice-wise model reduces the $\xHI$ value-distribution error by 65\%, the $b_1$ pull by 57\%, the bubble-size error by 54\% and small-scale incoherence by 45\%, and its fronts are 3.9\,Mpc wide against 6.2\,Mpc. During reionization, its $T_b$ distribution is 34\% closer to the truth. On the eight test cones inspected individually, it beats \texttt{fno} on every cone, by 12--28\% in full-cone RMSE. It also beats the 3-D rhizome C on every $\xHI$ morphology metric (value distribution $-39\%$, $b_1$ pull $-27\%$, bubble sizes $-19\%$, incoherence $-36\%$, edges $-19\%$), while $\tau_e$ and timing tie: with history channels, both rhizome models are already five to seven times better than \texttt{fno} there. Training on all cones instead of the 1,594-cone split improved per-cone RMSE by 3.6\% ($\xHI$) and 5.5\% ($T_b$) in the median, on 86\% and 75\% of cones respectively.

\begin{figure*}[!tp]
  \centering
  \figplaceholder[0.56\linewidth]{rhizome/h2h_cone24.png}{Cone 24 head-to-head.}
  \caption{Test cone 24. \emph{Top left:} $\xHI$ lightcone strips for truth, rhizome (slice-wise multi-field) and the \texttt{fno} baseline (labelled UFNO). \emph{Top right:} global history and signed error strips. \emph{Bottom:} transverse slices at mean neutral fraction 0.17, 0.40, 0.61 and 0.80, with per-slice RMSE. Full-cone RMSE 0.045 (rhizome) against 0.057 (\texttt{fno}). The \texttt{fno} fronts are blurred, and it leaves a faint haze inside ionized regions.}
  \label{fig:cone24}
\end{figure*}

\Cref{fig:h2h-ps} shows the scale dependence. Both models lose small-scale power late in reionization, \texttt{fno} more so. The rhizome keeps $r>0.9$ to higher $k$ at every stage, so the structure it does produce sits in the right place. \Cref{fig:ps-plane} puts this in the context of all models from the earlier 3-D operator study. In the plane of small-scale coherence against small-scale power error, the multi-field slice-wise rhizome sits apart from every \texttt{fno} model: $r(k>1)=0.90$ against 0.82--0.83, power error 0.17 against 0.31--0.34.

\begin{figure*}[!tp]
  \centering
  \figplaceholder[\linewidth]{rhizome/h2h_ps_stage_curves.png}{Power and coherence by stage.}
  \caption{$\xHI$ power ratio (left) and cross-correlation $r(k_\perp)$ (right) by reionization stage, multi-field slice-wise rhizome (top) and \texttt{fno} \texttt{tbs+histtb} (bottom), 199 test cones.}
  \label{fig:h2h-ps}
\end{figure*}

\begin{figure*}[!tp]
  \centering
  \figplaceholder[\linewidth]{rhizome/ps_plane_coherence_vs_power_best.png}{Coherence vs power plane.}
  \caption{Small-scale coherence $r(k>1\,\mathrm{Mpc}^{-1})$ against small-scale power error for the 3-D $\xHI$ models of the earlier operator study (open circles), the \texttt{fno} LOS-window models (filled) and the multi-field slice-wise rhizome (green), active stage. Better models lie to the lower right. Models from different pipelines and splits are compared with care.}
  \label{fig:ps-plane}
\end{figure*}

\begin{figure*}[!tp]
  \centering
  \figplaceholder[\linewidth]{rhizome/h2h_physical_history.png}{History and timing.}\\[4pt]
  \figplaceholder[\linewidth]{rhizome/h2h_bubble_size_summary.png}{Bubble size summary.}\\[4pt]
  \figplaceholder[\linewidth]{rhizome/h2h_edge_3d.png}{Edge profiles.}
  \caption{Physical metrics, rhizome (slice-wise multi-field) against \texttt{fno} \texttt{tbs+histtb}, 199 test cones. \emph{Top:} global history error, $\tau_e$ error in Planck $\sigma$ (median 0.01 against 0.08), and predicted against true midpoint and duration. \emph{Middle:} bubble-size distributions and their $W_1$ distance to the truth by stage (active 0.56 against 1.21\,Mpc). \emph{Bottom:} error profiles across ionization fronts; front width 3.9 against 6.2\,Mpc (truth 2.77).}
  \label{fig:h2h-physical}
\end{figure*}

\Cref{fig:h2h-errors} shows the signed error at $\bar x_{\mathrm{HI}}\approx0.4$ on eight test cones. As for the $\xHI$-only model, the rhizome's errors are thin rims on bubble walls, while \texttt{fno} misses whole regions and its rims are wider.

\begin{figure}[!htbp]
  \centering
  \figplaceholder[0.78\linewidth]{rhizome/h2h_error_maps_xhi04.png}{Error maps.}
  \caption{Signed $\xHI$ error maps at $\bar x_{\mathrm{HI}}\approx0.4$ on eight test cones (red: predicted too neutral), rhizome against \texttt{fno} \texttt{tbs+histtb}.}
  \label{fig:h2h-errors}
\end{figure}

\paragraph{Shared failure.} Some test cones contain large, flat plateaus of intermediate neutral fraction ($\xHI\approx0.3$--0.6) in the truth. No model reproduces them, and all overpredict the partially ionized fraction (0.13--0.18 against 0.09).

% ---------------------------------------------------------------------
\subsection{Brightness temperature before reionization}
\label{sec:results-absorption}

The slice-wise model has one clear weakness. Before reionization, when $T_b$ is in absorption, it produces too much small-scale $T_b$ power: $+24.8\sigma$ at $k=0.46\,\mathrm{Mpc}^{-1}$, against $+4.1\sigma$ for the 3-D rhizome and $+12.8\sigma$ for \texttt{fno} (\cref{tab:main-results}). The excess is concentrated in quiet slices whose true small-scale power is tiny; the power summed over the cone is in fact slightly low (ratio 0.97). The working explanation is that the spin variable $u$ is predicted slice by slice, and the per-pixel loss has no explicit term matching its fluctuation spectrum within a slice, whereas the 3-D model predicts it jointly along the LOS.

Three fixes that address data, post-processing and context did not work; a fourth, in the loss, did (\cref{tab:absorption-fixes}).

\begin{table}[!htbp]
  \centering
  \caption{Attempts to remove the absorption-era $T_b$ power excess of the slice-wise model (power deviation at $k=0.46\,\mathrm{Mpc}^{-1}$, plain model: $+24.8\sigma$).}
  \label{tab:absorption-fixes}
  \footnotesize
  \begin{tabular}{@{}p{2.2cm}p{1.05cm}p{4.3cm}@{}}
    \toprule
    Attempt & Excess & Outcome \\
    \midrule
    3.9$\times$ more training cones & $+24.8\sigma$ & $T_b$ values improve ($W_1$ $-31\%$), excess unchanged: no improvement from this data increase \\
    smooth $u$ along the LOS ($\sigma=4$ cells) & $+1.1\sigma$ & removes the excess only by destroying real LOS structure: $T_b$ RMSE $+40$--$110\%$, summed power $\times1.3$; rejected \\
    LOS context $\pm128$ cells & $+36.0\sigma$ & worse; per-voxel accuracy unchanged; not adopted \\
    \textbf{small-scale $T_b$ power loss} & $\bm{+1.5\sigma}$ & reduces the reported excess across scales; costs 4--8\% per-voxel RMSE \\
    \bottomrule
  \end{tabular}
\end{table}

The power loss compares the transverse power spectrum of each predicted slice with the truth in eight logarithmic bins from $k=0.3\,\mathrm{Mpc}^{-1}$ to the grid Nyquist frequency,
\begin{equation}
  \label{eq:power-loss}
  \mathcal{L}_{P} = \lambda\,\Big\langle \big(\log(P_{\mathrm{pred}}+\epsilon)-\log(P_{\mathrm{true}}+\epsilon)\big)^2 \Big\rangle_{\text{slices, bins}},
\end{equation}
with $\lambda=0.03$, applied to the mean-subtracted normalized $T_b$. It brings the absorption-era excess from $+24.8\sigma$ to $+1.5\sigma$, better than the 3-D rhizome, without overshooting the amplitude. It also improves the reionization-era statistics: $T_b$ power within 1--2\% at all scales, bubble-size $W_1$ $-14\%$ and $\tau_e$ $-20\%$ (\cref{fig:power-loss}). The price is a 4--8\% higher per-voxel RMSE, consistent over 80--90\% of the cones, and slightly worse edges. This reflects the different quantities optimized by the two losses: squared error rewards agreement at each cell, while the power term rewards agreement in fluctuation amplitude. Matching a power spectrum does not guarantee correct phases or morphology, so cross-correlation and front metrics remain necessary.

Applying the same term to $\xHI$ as well roughly halves the high-$k$ $\xHI$ power deficit (from 0.163 to 0.089 over active slices; in the latest stage from 0.59 to 0.15), gives the closest value distribution of all models ($W_1$ 0.0082), and keeps the absorption-era fix ($+1.3\sigma$). It costs another 2--4\% in per-voxel RMSE, and the extra small-scale power is only partly in phase with the truth (small-scale incoherence $+7\%$). The three variants are useful for different purposes: the plain model for per-voxel accuracy, the $T_b$ power loss for 21\,cm statistics, and both power terms for closer small-scale power amplitudes in both fields, with the stated phase-error trade-off.

\begin{figure*}[!tp]
  \centering
  \figplaceholder[\linewidth]{rhizome/power_loss_physical_21cm.png}{Power-loss 21cm statistics.}
  \caption{21\,cm statistics of the power-loss variants during reionization, 199 test cones: plain slice-wise model, $+\,T_b$ power loss (\texttt{pw0.03}), $+\,T_b$ and $\xHI$ power loss (\texttt{pw0.03\_px0.03}), and 3-D rhizome C. \emph{Right column:} power bias in units of single-slice sample variance for $T_b$ (top) and $\xHI$ (bottom). Both power-loss models stay near zero to $k\approx2\,\mathrm{Mpc}^{-1}$ in $T_b$; the plain model and 3-D C fall to $-10$ and $-12\sigma$.}
  \label{fig:power-loss}
\end{figure*}

A further third of the absorption-era $T_b$ error is a constant offset per slice, of random sign from cone to cone. Its size (1.6--2.0\,mK RMS) is nearly the same for all three model families and comparable to the error of the global-history emulator that provides the $\bar T_b(z)$ input (about 2.1\,mK). The similar error scales suggest that the history input may limit this component. Testing with true rather than emulated histories would separate that effect from model error. The power loss acts only on fluctuations and leaves the offset unchanged.

% ---------------------------------------------------------------------
\subsection{Controlled variations of the slice-wise model}
\label{sec:results-variations}

To separate components of the final model, the plain multi-field slice-wise configuration was retrained from scratch with one change at a time: a second seed, the number of updates, the interaction kernel, and the history inputs. Data, splits, losses, stopping rule and evaluation are the same as for the reference run (\cref{tab:variations}). The variants were trained with a later code version whose additions are inactive at these settings; the seed-1 run uses that same version. The kernel variants replace the global spectral kernel by a \emph{compact} real-space kernel supported on a disc of radius $R$ cells. All gates, the recurrence and the pointwise maps are kept. After $\Niter=6$ updates, information can then travel at most $6R$ cells. The rank-16 variant factorizes the channel mixing of the compact kernel.

\begin{table*}[!tp]
  \centering
  \caption{Controlled variations of the multi-field slice-wise model (width 64, all training cones, plain loss). The reference uses the spectral kernel with 24 modes, 6 updates and seed 0; each other row changes one setting and is retrained from scratch. Test RMSE on 3,129 held-out slices from 199 test cones; $\Delta$ relative to the seed-0 reference. Reach: largest distance information can travel in six updates. $P$ error: mean squared log ratio of the transverse power spectrum, predicted to true, in eight bins above $k=0.3\,\mathrm{Mpc}^{-1}$ (not recorded for the reference run). Training time on one A100, including validation.}
  \label{tab:variations}
  \footnotesize
  \setlength{\tabcolsep}{4pt}
  \begin{tabular}{@{}lrrrrrr@{}}
    \toprule
    Variant & Params & Reach & Time (h) & $\xHI$ RMSE ($\Delta$) & $T_b$ RMSE, mK ($\Delta$) & $P$ error $T_b$ / $\xHI$ \\
    \midrule
    reference (seed 0) & 4.57\,M & global & 3.6 & 0.0872 & 7.71 & -- \\
    second seed (seed 1) & 4.57\,M & global & 3.1 & 0.0873 ($+0.2\%$) & 7.82 ($+1.3\%$) & 0.385 / 0.392 \\
    3 updates & 4.57\,M & global & 2.2 & 0.0866 ($-0.6\%$) & 7.68 ($-0.5\%$) & 0.395 / 0.385 \\
    10 updates & 4.57\,M & global & 4.9 & 0.0869 ($-0.3\%$) & 7.66 ($-0.7\%$) & 0.378 / 0.355 \\
    \midrule
    compact kernel, $R=1$ & 66\,k & 8.6\,Mpc & 2.8 & 0.1209 ($+38.7\%$) & 9.76 ($+26.5\%$) & 0.511 / 0.449 \\
    compact kernel, $R=2$ & 99\,k & 17\,Mpc & 3.8 & 0.1089 ($+24.9\%$) & 8.85 ($+14.7\%$) & 0.424 / 0.398 \\
    compact kernel, $R=3$ & 164\,k & 26\,Mpc & 4.4 & 0.1010 ($+15.9\%$) & 8.27 ($+7.3\%$) & 0.424 / 0.402 \\
    compact kernel, $R=4$ & 246\,k & 34\,Mpc & 5.2 & 0.0972 ($+11.5\%$) & 8.10 ($+5.1\%$) & 0.417 / 0.385 \\
    compact kernel, $R=6$ & 508\,k & 51\,Mpc & 7.7 & 0.0951 ($+9.1\%$) & 8.13 ($+5.4\%$) & 0.444 / 0.437 \\
    compact kernel, $R=8$, rank 16 & 51\,k & 69\,Mpc & 4.6 & 0.0926 ($+6.3\%$) & 8.04 ($+4.2\%$) & 0.444 / 0.434 \\
    \midrule
    no $\bar x_{\mathrm{HI}}(z)$ history input & 4.57\,M & global & 3.7 & 0.0973 ($+11.7\%$) & 8.26 ($+7.1\%$) & 0.399 / 0.431 \\
    no $\bar T_b(z)$ history input & 4.57\,M & global & 3.9 & 0.0900 ($+3.3\%$) & 8.04 ($+4.3\%$) & 0.341 / 0.400 \\
    \bottomrule
  \end{tabular}
\end{table*}

\paragraph{Seed and rollout length.} The second seed differs from the reference by 0.2\% in $\xHI$ and 1.3\% in $T_b$ RMSE. This is the only direct estimate of training noise in this report, and it suggests that differences of one or two percent between single runs should not be interpreted. Three and ten updates match six within that spread. Three updates train in 2.2\,h instead of 3.6\,h. Between 3 and 10 updates, the rollout length is therefore not a sensitive setting on this task. This differs from the toy pilot of \cref{tab:rhizome-pilot}, where 1 or 2 updates were clearly worse, but there each update count was evaluated without retraining.

\paragraph{Global reach of the kernel.} Every compact kernel is worse than the spectral kernel, and the error falls steadily as the reach grows: from $+39\%$ in $\xHI$ at 8.6\,Mpc to $+6\%$ at 69\,Mpc. The compact kernels also have 9--90 times fewer parameters, so reach and capacity change together. Within the compact family, however, the error follows reach rather than parameter count: the rank-16 kernel with $R=8$ has the fewest parameters (51\,k) and the smallest error, while $R=6$ with ten times more parameters is worse. The compact kernels are not cheaper to train on an A100 either, because wide direct convolutions are slower than the FFT (7.7\,h at $R=6$). These runs support the global-reach design principle of \cref{sec:rhizome-principles}. Since every compact variant keeps the gates, they do not measure the value of the gates.

\paragraph{Global-history inputs.} Both histories were removed in turn from the scalar inputs. Paired over the 199 full test cones, removing the $\xHI$ history raises the median per-cone RMSE by 12.6\% in $\xHI$ (95\% bootstrap interval 11.5--14.5\%) and by 11.3\% in $T_b$; it is worse on 91\% and 93\% of the cones. The error of the global history $\bar x_{\mathrm{HI}}(z)$ grows by 38\%. Removing the $T_b$ history raises the median per-cone $T_b$ error by 26.6\% (interval 19.7--31.2\%), and by 46\% before reionization. Its $\xHI$ error rises by 6.1\%. Both changes are far larger than the seed spread. They agree with the 3-D ablation of \cref{tab:results-3d-mf}: the emulated histories carry information that the local density bands do not. One effect goes the other way: without the $T_b$ history the slice-level $T_b$ power error is lower (0.341 against 0.378--0.395), so the history input improves per-voxel accuracy but not the small-scale fluctuation amplitude.

% ---------------------------------------------------------------------
\subsection{Learned spectral weights}
\label{sec:results-modes}

Neural networks tend to learn low frequencies first, and a Fourier layer can end up using only its lowest modes. \Cref{fig:mode-weights} tests this for the rhizome by comparing the RMS weight of each Fourier mode of the trained kernel with its value at initialization. The weights grew 13--45-fold almost evenly up to the band edge: the ratio of high-quarter to low-quarter weight is 0.78--0.94 for all rhizome models, against 0.10--0.45 for a U-FNO in an earlier study. Models with 36 or 48 modes also develop non-negligible weights in their extra modes. Weight amplitude alone, however, does not measure mode usage: the transmitted input spectrum and output sensitivity also matter. These profiles show no collapse of the multiplier weights, but do not establish why the 48-mode sweep performed worse or rule out an alternative kernel parameterization. Longer matched-budget runs and per-mode message or sensitivity measurements would test those explanations.

\begin{figure*}[!tp]
  \centering
  \figplaceholder[0.82\linewidth]{rhizome/mode_weight_profiles.png}{Mode weight profiles.}
  \caption{Per-mode RMS of the learned Fourier kernel relative to initialization. Radial profiles over wavenumber; they are flat to the band edge: there is no collapse onto low frequencies.}
  \label{fig:mode-weights}
\end{figure*}

% ---------------------------------------------------------------------
\subsection{Learned gates on 21\,cm slices}
\label{sec:results-gates}

\Cref{fig:mechanism} inspects the gates of a toy model. The same analysis was repeated on the trained multi-field slice-wise model, averaged over 120 test slices with $0.1\le\bar x_{\mathrm{HI}}\le0.9$, one per cone. Receiver cells are classified with the true field: 3,808 cells on ionization fronts, and 1,920 and 1,489 cells deep inside neutral and ionized regions. For each receiver $i$ and update $t$, the re-weighting $M_t(j)$ of \cref{fig:mechanism} is averaged radially in $|r_i-r_j|$ (\cref{fig:mf-gates}).

\begin{figure*}[!tp]
  \centering
  \figplaceholder[0.82\linewidth]{rhizome/mf_gate_gain_receiver.png}{Gate gain and receiver gate per update.}
  \caption{Gates of the trained multi-field slice-wise model on 120 test slices, by receiver class and update. \emph{Left:} mean gain $\langle M_t\rangle_j$ of the incoming connections. \emph{Right:} channel-RMS receiver gate $\norm{A_i}_{\mathrm{rms}}$. Solid: plain model; dashed: model trained with both power terms (\cref{sec:results-absorption}). Gains on the transmitted value, not output sensitivities.}
  \label{fig:mf-gates}
\end{figure*}

\paragraph{The gates depend on the state.} The receiver gate separates the phases, and the separation grows during the rollout. In the plain model, the channel-RMS receiver gate of ionized-interior cells rises from 0.34 to 0.46 over the six updates, while that of neutral-interior cells falls from 0.26 to 0.15 and that of front cells from 0.25 to 0.17. At the last update ionized cells accept about three times more of the incoming message than neutral ones. The source gates stay close to one (channel RMS 0.81--0.97 on four inspected slices), so, as in the toy model, the receiver dominates. The gates do not merely rescale: the radial profile of the re-weighting also changes. For front receivers, the boosted region contracts from a half-width of 26\,Mpc at the first update to 14\,Mpc at the last, with a peak boost of $2.0\times$ at 3.6\,Mpc; for neutral-interior receivers it contracts from 28 to 5\,Mpc. For ionized-interior receivers it stays broad (25 to 24\,Mpc). Beyond 40\,Mpc the re-weighting is flat for all classes. On four single slices, the normalized shape of the effective connection map into a front cell changes by 25--30\% between the first and last update, against about 4\% in the toy model. The model trained with both power terms has larger receiver gates and concentrates its re-weighting more strongly, within about 1\,Mpc of the receiver ($3.5\times$ at fronts).

\paragraph{What this does and does not show.} The trained gates are clearly state dependent: they distinguish ionized from neutral cells, and the connections they produce change over the rollout, mostly by narrowing towards nearby sources in neutral and front cells. This is a description of the learned mechanism, not a measurement of its value. A model with gates computed once from the initial state, or without gates, could reach the same accuracy by other means, and $M_t$ measures gains on the transmitted value rather than influence on the output.

% ---------------------------------------------------------------------
\subsection{Cost}
\label{sec:results-cost}

Per training slice, the 2-D model is about four times slower than the 3-D models, since it runs full-resolution recurrent updates in float32. It needs 15--20 times fewer training slices, however (2.2\,M against 33--49\,M), so a complete training run is four to seven times cheaper: 3.5\,h on one A100 for the best model, against about 18\,h on an H200 for the 3-D rhizome and 44.7\,h on an A100 for \texttt{cnn\_whno\_es}.

At inference, a full multi-field cone takes 5.5\,s end to end on an A100 and 3.7\,s on an H200 once the LOS bands are built on the GPU and the next data window is prefetched (\cref{fig:inference-cost}). This makes the slice-wise model the fastest model end to end, 1.8--2.3 times faster than both 3-D models; at that point inference is limited by reading the data. The GPU forward pass alone (3.8\,s / 1.7\,s) matches the 3-D rhizome and is about twice as slow as \texttt{fno}. Mixed precision is not worthwhile: bfloat16 gains only 5--14\%, because the FFTs stay in float32, and it shifts $\xHI$ by up to 0.08.

\begin{figure}[!htbp]
  \centering
  \figplaceholder[\linewidth]{inference_cost.pdf}{Inference cost chart.}
  \caption{End-to-end inference time for one full multi-field test cone ($\xHI$ and $T_b$), including data reading.}
  \label{fig:inference-cost}
\end{figure}

% ---------------------------------------------------------------------
\subsection{Ablations}
\label{sec:ablations}

\Cref{tab:ablations} collects what each design and training choice was worth across the experiments above. The table includes controlled changes as well as comparisons of model configurations. All training runs use a single seed, except one seed repeat of the main configuration, which differs by 0.2\% ($\xHI$) and 1.3\% ($T_b$) in test RMSE. The 2-D/3-D comparison changes several architectural and numerical choices together.

\begin{table*}[!tp]
  \centering
  \caption{Summary of ablations. Effects are relative changes in the stated metric; negative is better.}
  \label{tab:ablations}
  \small
  \begin{tabular}{@{}p{4.6cm}p{3.6cm}p{5.6cm}@{}}
    \toprule
    Change & Metric & Effect \\
    \midrule
    LOS bands vs single slice (2-D) & mixed-slice RMSE & $-36\%$ (0.174 $\to$ 0.112) \\
    width 64 $\to$ 128 (2-D) & mixed-slice RMSE & $-3\%$ at $4\times$ parameters \\
    width 48 $\to$ 64 (3-D) & test RMSE & $+6\%$ (worse) \\
    lr $2\times10^{-3}\to3\times10^{-4}$ (3-D) & test RMSE & $+13\%$ (worse) \\
    warm restart (3-D) & test RMSE & $-3$ to $-6\%$ \\
    history channels (3-D, A $\to$ B) & RMSE $\xHI$ / $T_b$ & $-21\%$ / $-34\%$ \\
    structured $T_b$ head (B $\to$ C) & $T_b$ RMSE; timing & $-5\%$; $-15$ to $-30\%$ \\
    slice-wise 2-D vs 3-D (same data) & $1-r$; edge $H^1$ & $-40\%$; $-21\%$ \\
    all 6,169 vs 1,594 training cones & median per-cone RMSE & $-3.6\%$ ($\xHI$), $-5.5\%$ ($T_b$) \\
    LOS context $\pm128$ vs $\pm32$ & per-voxel RMSE & little reported change; duration $-37\%$; absorption $T_b$ worse \\
    $T_b$ power loss & pre-reion.\ $T_b$ power dev. & $+24.8\sigma\to+1.5\sigma$; RMSE $+4$ to $+8\%$ \\
    $+\,\xHI$ power loss & high-$k$ $\xHI$ power error & $-45\%$; RMSE further $+2$ to $+4\%$ \\
    \midrule
    \multicolumn{3}{@{}l}{\emph{Slice-wise multi-field model, retrained (\cref{sec:results-variations})}} \\
    second seed & test RMSE $\xHI$ / $T_b$ & $+0.2\%$ / $+1.3\%$ (seed spread) \\
    3 or 10 updates instead of 6 & test RMSE $\xHI$ / $T_b$ & $-0.3$ to $-0.6\%$ / $-0.5$ to $-0.7\%$ (within seed spread) \\
    compact kernel, reach 8.6--69\,Mpc, vs spectral & test RMSE $\xHI$ / $T_b$ & $+39\%$ to $+6\%$ / $+27\%$ to $+4\%$ (worse; shrinks with reach) \\
    remove $\bar x_{\mathrm{HI}}(z)$ history & median per-cone RMSE $\xHI$ / $T_b$ & $+12.6\%$ / $+11.3\%$ (worse) \\
    remove $\bar T_b(z)$ history & median per-cone RMSE $\xHI$ / $T_b$ & $+6.1\%$ / $+26.6\%$ (worse) \\
    \bottomrule
  \end{tabular}
\end{table*}

\paragraph{What the variations establish about the operator.} The retrained variations of \cref{sec:results-variations} isolate three components of the operator. The global reach of the kernel matters: no compact kernel up to 69\,Mpc reach matches the spectral kernel. The rollout length matters little between 3 and 10 updates. The history inputs carry information the operator cannot recover from the density bands. \Cref{sec:results-gates} shows that the trained gates are state dependent and change over the rollout.

\paragraph{What remains open.} None of these runs removes or freezes the gates, so the accuracy contribution of the state-dependent gating is still not measured. The planned study compares the full rhizome with six otherwise identical arms (no gates, source-only, receiver-only, gates computed from the conditioning only, gates computed once from the initial state, and static per-channel gates), five seeds each, all retrained from scratch. With the seed spread now measured at 1--2\%, the primary contrast, full against initial-state gates, needs differences well above that level to be interpretable. Until then, the comparison with \texttt{fno} establishes that the rhizome \emph{as a whole} works better on this task, and the variations attribute part of that to the global kernel and the inputs, but not to the state-dependent gating specifically.

% ---------------------------------------------------------------------
\subsection{Discussion}
\label{sec:discussion}

\paragraph{What the rhizome does well.} Across every configuration, the rhizome's advantage over the \texttt{fno} models is larger on physical observables than on per-voxel RMSE. It places ionization fronts more accurately (higher small-scale coherence), keeps them sharper, and reproduces bubble-size distributions, bias and stochasticity more faithfully. These are among the quantities that \cref{sec:fno-limitations} identified as useful diagnostics of a restricted spectral interaction. They are also the ones that matter for 21\,cm inference, where bubble morphology and the field-level bias enter the observables.

\paragraph{What the line of sight needs.} The most consequential finding is not about the operator itself. A 2-D model with $\pm32$ cells of LOS context, applied slice by slice, outperforms native 3-D models trained on the same data, and does so at a fraction of the cost. Together with the periodic tiling of the lightcones, this suggests that the bottleneck of the 3-D models is how they handle the LOS, not their transverse expressivity.

\paragraph{Scope of the evidence.} The targets are semi-numerical lightcones from one simulation suite with repeated periodic boxes. Accuracy on those targets does not by itself establish agreement with resolved radiative transfer, unseen cosmologies, another box size or observational inference. The emulator is deterministic and has no calibrated field uncertainty. Power spectra, morphology diagnostics and optical-depth errors assess different aspects of agreement, but do not replace a downstream inference test.

\paragraph{Reproducibility details still needed.} The optimizer, stopping rule, normalization, loss weights and slice sampling of the slice-wise models are now recorded in \cref{sec:training}. Before treating this report as a complete methods paper, the simulation-parameter ranges and sampling law, split identifiers, the corresponding settings of the 3-D models, history-emulator training split, and precise metric definitions should also be recorded. The history inputs need an explicit provenance check showing that test-cone targets were not used to fit or select the emulator. LOS padding length must be distinguished from core and halo lengths. The meaning of the quoted $\sigma$ and $b_1$ pulls, the partially ionized thresholds, and power-bin boundaries also need to be specified. These are missing documentation or checks, not assumptions supplied by this revision.

\TODO{Record whether the test comparisons informed the choice of context, loss weights or checkpoints.}

\paragraph{Caveats.} Several confounds remain. Nearly all results are single seeds. A repeat of the main slice-wise configuration differs by 0.2--1.3\% in test RMSE, so differences of a few percent (for example A against A0, which is a 3-D run with its own noise level) may be partly noise. Part of the advantage in $\tau_e$ and the global history comes from training with BCE rather than from the operator, and the \texttt{fno} models should be retrained with BCE to separate the two. The contribution of the gating itself has not been isolated (\cref{sec:ablations}), although the trained gates are demonstrably state dependent (\cref{sec:results-gates}). Finally, all results are at one resolution; the grid-independence of the parameters has not been put to the test.


